% Chapter 8 -- affinity. Update with affinity/* (demand.py, engine.py, weights.py, calibration.py, knowledge.py,
% preferences.py, batch.py, texts.py, data/directions.json, data/knowledge.json), gui/affinity_box.py,
% gui/affinity_service.py, gui/affinity_settings.py, gui/window_affinity.py.
\chapter{Affinity: the score and the verdict}\label{ch:affinity}

Chapter~\ref{ch:scoring} scores a company against what you \emph{prefer} (distance, contract) and
a posting against the CV. The \textbf{affinity} answers the question of the product: \emph{how
well does this profile fit what this company seems to need?} It is a number from 0 to 100\,\%, a
confidence, a \textbf{verdict} and the reasons behind them. It is computed locally, deterministically, from a CV and the data on the company card, and is always
an estimate, never a prediction of being hired. Company-type preferences only move a company
in the lists; whether it is hiring or growing is a separate \emph{opportunity} signal, not part of the number.

\section{What the company asks for: the demand}
The \emph{demand} is assembled from the card (or a search result) as a list of items, each with an
origin and a weight. Postings are read as in chapter~\ref{ch:scoring}; everything else is read from
the company's own texts, and an item already present with a stronger origin is not repeated.

{\small
\begin{tabularx}{\textwidth}{@{}l c Y@{}}
\toprule
\hdr \thead{Origin} & \thead{Weight} & \thead{Where it comes from} \\
\midrule
Must (explicit)    & 5 & A requirement line of a posting classed ``must''. \\
Nice (explicit)    & 3 & A ``nice to have'' line. \\
Context (explicit) & 1 (+1 for a duty) & Other posting text; counts only as positive evidence. \\
Stated             & 2 & A technology from the skills dictionary named on the company's site, products or presentation. \\
Inferred from products & 0.3--0.7 & Product rules of the knowledge table (section~\ref{sec:knowledge}), e.g.\ packaging machines suggest PLC and HMI skills. \\
Inferred from the type & 0.3--0.6, $\times$0.8 or 0.5 & What a kind of company usually employs (machine builder, system integrator, manufacturer, software house). \\
\bottomrule
\end{tabularx}
}

\vspace{3pt}
Besides skills the demand lists \emph{roles} (posting titles 1.0, inferred from products 0.8, from the
type 0.7), \emph{sectors} (ATECO 1.0, OpenStreetMap 0.8, site keywords 0.6) and \emph{asked}
values (years, languages, degree, seniority words in a title). An \emph{inferred} item is a guess, labelled ``inferred from their products
(not a requirement)''; the company name classifies the type only, never counts as a skill.

\section{Directions and components}
A \textbf{direction} is a kind of career defined by role words and skill families. Four ship enabled:
\emph{Industrial automation}, \emph{Software development}, \emph{Cybersecurity and industrial
networks} and \emph{Technical coordination}; Settings, \ui{Affinity}, switches them and adds your own. The company is scored once per applicable direction, using only the
demand that belongs to it, and once on the whole profile. A direction is \emph{applicable} when the company has a posting in it or at least 15\,\% of its demand lies in it; otherwise it is not scored (no penalty). The headline is the best applicable direction, shown as ``Fits you as:
<direction>, 78\,\% (confidence)'' with the others as ``Also as\ldots''; with no direction applicable
it is the whole profile. A direction whose best posting fits gets the same bonus as Apply in the lists, so it wins the headline over a direction that scores a little higher without an open posting.

Each score is a weighted mean of six components (0--100); one without data is dropped and the weights renormalised.

{\small
\begin{tabularx}{\textwidth}{@{}l c Y@{}}
\toprule
\hdr \thead{Component} & \thead{Weight} & \thead{How it is scored} \\
\midrule
Skills       & 0.35 & Each demand item gets the CV evidence of chapter~\ref{ch:analysis}: named in a dated job 1.0, elsewhere 0.85, adjacent 0.5, absent 0; times recency (1 up to 24 months, decaying to 0.6; only listed 0.8). With a posting: 70\,\% best-posting coverage plus 30\,\% overlap with the stated and inferred skills. \\
Role         & 0.20 & Company roles against your current and past titles, Italian and English variants (fuzzy and by meaning); a role inferred from the type can reach 70 at most. \\
Sector       & 0.15 & Closeness of the company's sector to the ones of your jobs (months, recency) and to your ATECO categories. \\
Seniority    & 0.10 & Years asked against years held (too few: held/asked; far too many: down to 0.5) and level words; only with a posting that asks. \\
Constraints  & 0.10 & The knock-out answers of chapter~\ref{ch:analysis} (ok 1, unknown 0.8, warn 0.6, fail 0) and the distance limit if set. \\
Practical    & 0.10 & Contract, remote work and salary against your preferences, where the posting publishes them. \\
\bottomrule
\end{tabularx}
}

\section{Hard requirements, coverage, confidence}
\textbf{Hard-requirement cap.} If a screening question of the best posting is \emph{failed} (a must
for a language, degree, licence, car, travel, right to work\ldots\ that your saved answers contradict),
if the posting is gated (chapter~\ref{ch:scoring}), or if the company is beyond your distance limit,
the affinity is capped at \textbf{25} and the verdict is \emph{Not a fit}.


\textbf{Postings known only by title.} A posting read as a bare title (a careers page listing ``Back-end developer'') is not scored 0: the title is compared with your CV roles. An internship, apprenticeship or graduate programme (intern, stage, tirocinio, apprendistato, trainee\ldots) never counts as fitting a CV with five or more years of experience, and the card says why.

\textbf{Coverage and ``not enough data''.} Each component counts only as far as its evidence is firm
(a posting item 1.0, a skill stated on the site 0.8, inferred from products 0.5, from the type only 0.25).
\emph{Coverage} is the share of the possible weight that rests on evidence. Below \textbf{0.45}, or
with neither skills nor role known, there is no number: the card says \emph{Not enough data} and what
to collect (add the website, build the full card, paste a posting, add the ATECO code).

\textbf{Confidence}: \emph{high} from coverage 0.75 with a described posting, \emph{medium} from 0.55, else \emph{low}.

\section{The verdict}\label{sec:order}
The first rule that applies wins; the thresholds are those of Settings after calibration.

{\small
\begin{tabularx}{\textwidth}{@{}l Y@{}}
\toprule
\hdr \thead{Verdict} & \thead{When} \\
\midrule
Not enough data & No affinity number (above). \\
Not a fit       & Company type switched off (0\,\%), a hard requirement failed, ``directions only'' is on and the company is outside every direction, or affinity below \textbf{40}. \\
Apply           & The best posting fits (its match score is at least 55, not capped, no failed requirement) and affinity is at least \textbf{50}: send your CV. \\
Spontaneous application & Affinity at least \textbf{55}, a \emph{plausible need} (stated or inferred skills overlap at least 40\,\% with one item at 50 or more, or the role scores 50 or more) and confidence not low. \\
Keep an eye on it & Everything else: it fits, but nothing is open and the signals are weak. \\
\bottomrule
\end{tabularx}
}

\vspace{3pt}
\textbf{Ordering.} Lists sort by tier: \emph{Apply} and \emph{Spontaneous} first, then
\emph{Keep an eye on it}, \emph{Not a fit}, \emph{Not enough data}. Inside a tier the key is the affinity
plus 15 for \emph{Apply} (postings first), times your company-type weight; ties go to higher
confidence, then the name. So an \emph{Apply} at 60\,\% (key 75) beats a \emph{Spontaneous} at 74\,\%; a lower tier never outranks the first. A column header overrides this.

\section{Calibration from your judgements}
In the box on the card, \ui{Agree} and \ui{Disagree} record your opinion of one result (the last one
counts). From \textbf{8} judgements Settings, \ui{Affinity}, offers \ui{Calibrate from my judgements}:
a \emph{disagreement} with a positive verdict pushes the component weights and the thresholds towards
fewer positives, one with \emph{Not a fit} the other way. Disagreements below 20\,\% of the judgements
change nothing; a weight moves at most 0.05 a step (staying within 0.03--0.45), thresholds by 2 points, never more
than 10 from their defaults (40, 50, 55). \ui{Apply these weights} saves
them and recomputes, \ui{Default weights} resets. Nothing else is calibrated.

\section{The conservative knowledge table}\label{sec:knowledge}
For a company with no posting the demand is \emph{inferred} from a small hand-made, editable table
(\code{affinity/\allowbreak data/\allowbreak knowledge.json}): 14 product rules (packaging, automatic machines, robotic cells,
control panels, supervision and MES, machine tools, material handling, business and custom software,
cybersecurity services, IT infrastructure, OT security, engineering services, production plants)
with a few skills and roles each, and four company types with text, ATECO and map-tag evidence. All
weights are at most 0.7, type-derived ones are cut again, inferred items are labelled, and a result resting on the type
alone has coverage too low to give a number. A table of sector closeness completes it.

\section{What it cannot know}
Only what is on the card and in your CV text: a skill you have but did not write counts as a gap (``if you have
it, state it''). Meaning-based matching of titles needs the language model already on disk. It does not know the
employer's real needs, the competition, or whether a post is filled.

\section{Where you see it}
The \textbf{box at the top of the company card} shows the verdict badge, the percentage with its
confidence, the best direction and company type, and one line of reasons; \ui{Details} lists per component the evidence with source links, strengths and gaps, up to three \emph{points for
the cover letter} (\ui{Into the e-mail} puts them in the draft, chapter~\ref{ch:apply}) and ``What to collect''. The columns \ui{Affinity}, \ui{Verdict}, \ui{Fits as} and \ui{Type} appear in every company list. Results are cached and recomputed when the CV, the company data, the settings or the month
change. Settings, \ui{Affinity}: directions, kinds of company (0--100\,\%; 0 excludes), an optional
distance limit (off; 1--2\,000\,km) and calibration.
